\documentclass[11pt,reqno]{amsart}
\usepackage[T1]{fontenc}
\usepackage{amsmath,amssymb,amsthm,mathtools,array,longtable}
\usepackage[margin=1in]{geometry}
\usepackage[colorlinks=true,linkcolor=blue,urlcolor=blue]{hyperref}
\newtheorem{theorem}{Theorem}[section]
\newtheorem{lemma}[theorem]{Lemma}
\newtheorem{proposition}[theorem]{Proposition}
\newtheorem{corollary}[theorem]{Corollary}
\theoremstyle{definition}
\newtheorem{definition}[theorem]{Definition}
\newtheorem{remark}[theorem]{Remark}
\newcommand{\R}{\mathbb R}
\newcommand{\C}{\mathbb C}
\newcommand{\Z}{\mathbb Z}
\newcommand{\N}{\mathbb N}
\newcommand{\E}{\mathbb E}
\newcommand{\Sph}{\mathbb S^2}
\newcommand{\dd}{\mathrm d}
\newcommand{\Var}{\operatorname{Var}}
\newcommand{\gap}{\operatorname{gap}}
\newcommand{\TV}{d_{\mathrm{TV}}}
\newcommand{\dist}{\operatorname{dist}_\infty}
\newcommand{\Ran}{\operatorname{Ran}}
\newcommand{\af}{\mathrm{af}}
\newcommand{\ff}{\mathrm f}
\newcommand{\cD}{\mathcal D}
\newcommand{\Lip}{\operatorname{Lip}}
\title{The spin-2 square-lattice AKLT gap on tori of every parity}
\author{STRENGTHEN-ALL: replacement and added manuscript text}
\date{November 16, 2026}
\begin{document}
\begin{abstract}
The evenness assumption on both torus periods can be removed. Three
explicit, disjoint, antipodal spherical designs give positive classical
weights on every torus. The sign change is made only after lifting a
conditional box to the plane. The existing continuous estimate then
gives a uniform window bound for the resulting inhomogeneous finite-state
specifications. A finite-state box-gap argument requires no translation
invariance. Uniqueness holds for every pair of periods at least three;
more generally the ground space on every finite simple graph of maximum
degree four is computed exactly.
\end{abstract}
\maketitle

\section{Replacement for the statement and proof structure in Section 1}
\label{s:intro}

We use the manuscript's spin-2 polynomial Hilbert space, edge projections
and normalization of the Hamiltonian. In particular a bond is an element
of a simple edge set, and the heat-bath Dirichlet form has one update per
site. For $L,N\ge3$, let $E$ be the nearest-neighbour edge set of
$\Lambda=\Z_L\times\Z_N$, and fix an arbitrary orientation of every edge.
Write
\[
 q_{xy}=z_{x,1}z_{y,2}-z_{x,2}z_{y,1},\qquad
 B=\prod_{e\in E}q_e,\qquad H_{L,N}=\sum_{e\in E}P_e.
\]

\begin{theorem}[Replacement for Theorem 1.1]\label{t:main}
There are $c>0$ and $L_0\in\N$ such that, for every pair of integers
$L,N\ge L_0$, with no parity restriction,
\[
 \ker H_{L,N}=\C B,\qquad \gap(H_{L,N})\ge c.
\]
In addition, $\ker H_{L,N}=\C B$ for every $L,N\ge3$.
\end{theorem}

The uniform-gap assertion uses the same computer-assisted inputs as the
manuscript, at $(c_{\mathrm s},c_{\mathrm t})=(961/1000,187/200)$.
Throughout this text the unchanged analytic statements are understood with
the adjudicated corrections in ADJ-R4, ADJ-R5, ADJ-R6 and ADJ-MINOR.
Their corrected implication (H1), (H2), (H3$'$)$\Rightarrow$(P) is retained.
The present text proves every additional implication needed to remove
parity. It does not assert an independent replay of the large certificates.

Replace the first three reductions in the introduction as follows.
Assign one of three rotated icosahedral frames according to a proper
three-colouring of the torus. The squared modulus of the frame image of
$B$, multiplied by the reference weights and normalized, is a strictly
positive antiferromagnetic measure. Its heat-bath gap bounds the quantum
gap below. The classical finite-size criterion is already valid without
parity or translation invariance. A conditional box lifts to the plane,
where a sitewise sign change makes every interaction ferromagnetic.
The resulting alphabets depend on the sites. Uniform quadrature and the
continuous cut argument give the window estimate at every translate;
the finite-state proof of MOSPROOF then gives uniform box gaps without
translation invariance. There is no global sign change on an odd torus.

The uniqueness assertion is independent of all mixing estimates and
computer-assisted inputs, by Proposition~\ref{p:graphkernel} below.

\section{Replacement and additions in Section 2}
\label{s:frames}

\subsection{Ground spaces on simple graphs}

The manuscript's edge lemma remains unchanged: in the polynomial
realization $V=\operatorname{Sym}^4(\C^2)$,
$P_{\{x,y\}}\varphi=0$ if and only if $q_{xy}$ divides $\varphi$.
For completeness, multiplication by $q_{xy}$ embeds the bidegree-$(3,3)$
polynomials equivariantly into the bidegree-$(4,4)$ polynomials. By
Clebsch--Gordan its image consists precisely of spins $0,1,2,3$, the
orthogonal complement of spin $4$. Tensoring with the other sites proves
the assertion for arbitrary $\varphi$.

\begin{proposition}[Exact ground space]\label{p:graphkernel}
Let $G=(V_G,E_G)$ be a finite simple graph with $d_x=\deg_G(x)\le4$.
Put $B_G=\prod_{e\in E_G}q_e$, and let $\mathcal P_{(4-d_x)}$ denote
the polynomials homogeneous of degree $4-d_x$ in each $z_x$.
Multiplication by $B_G$ is a linear isomorphism
\begin{equation}\label{e:graphkernel}
 \mathcal P_{(4-d_x)}\ \longrightarrow\ \ker\Bigl(\sum_{e\in E_G}P_e\Bigr),
 \qquad \chi\longmapsto B_G\chi.
\end{equation}
Consequently
\[
 \dim\ker H_G=\prod_{x\in V_G}(5-d_x).
\]
Every finite simple $4$-regular graph has the unique ground state $B_G$,
up to a scalar, even when it is disconnected or nonbipartite.
\end{proposition}
\begin{proof}
Each determinant $q_e$ is irreducible in $\C[z_{V_G}]$: as a quadratic
form in its four endpoint variables it has rank four, whereas a product
of two linear forms has rank at most two. Irreducibility persists after
adjoining the other variables, since the quotient by $q_e$ remains an
integral domain. Distinct edges give nonassociate irreducibles because
their endpoint variable sets differ. In the polynomial unique
factorization domain, if every $q_e$ divides $\varphi$, their product
divides $\varphi$. Positivity of the projections gives
$\ker\sum_eP_e=\bigcap_e\ker P_e$, so the edge lemma proves that every
ground state is $B_G\chi$. Multihomogeneity forces $\chi$ to have degrees
exactly $4-d_x$: decompose $\chi$ into its multihomogeneous components
and use that multiplication by the nonzero multihomogeneous polynomial
$B_G$ cannot annihilate any component. Conversely every such product
is killed by each $P_e$. Multiplication is injective in an integral
domain. The residual polynomial space has dimension
$\prod_x(4-d_x+1)$, proving all assertions.
\end{proof}

For $L,N\ge3$ the four neighbours $(i\pm1,j)$ and $(i,j\pm1)$ are
distinct and differ from $(i,j)$. Thus the torus is simple and
$4$-regular, also when a period is three. Proposition~\ref{p:graphkernel}
proves uniqueness for all sizes under consideration. Reversing an edge
orientation multiplies $B$ by $-1$ and does not change its span.

\subsection{Three exact frames and their separation}

Keep the definition of a positive weighted spherical $4$-design and the
tight-frame lemma of the manuscript. Replace the checkerboard construction
and its icosahedron--dodecahedron bound by the following construction.
Let
\begin{equation}\label{e:threeframes}
 \begin{gathered}
 \phi=\frac{1+\sqrt5}{2},\quad r^2=\phi+2,\quad
 D=\frac1r\{(0,\pm1,\pm\phi),(\pm1,\pm\phi,0),(\pm\phi,0,\pm1)\},\\
 T=\begin{pmatrix}99/101&-20/101&0\\20/101&99/101&0\\0&0&1\end{pmatrix},
 \qquad D^{(j)}=T^jD\quad(0\le j\le2).
 \end{gathered}
\end{equation}
All signs in the displayed vertices are independent; all frames have
uniform weights $1/12$.

\begin{lemma}[Designs and exact minimum weight]\label{l:threeframes}
Each $D^{(j)}$ is an antipodal spherical $5$-design. Set
\begin{equation}\label{e:delta}
 b=\frac{5-\sqrt5}{505}=\frac{2}{101(\phi+2)},\qquad \Gamma=1-b.
\end{equation}
For distinct $i,j$ and $a\in D^{(i)},a'\in D^{(j)}$,
\begin{equation}\label{e:separation}
 |a\cdot a'|\le\Gamma<1,\qquad b\le1\pm a\cdot a'\le2-b.
\end{equation}
The lower bound $b$ is attained. In particular the three designs are
pairwise disjoint, including up to an antipodal sign.
\end{lemma}
\begin{proof}
The independent coordinate sign changes in $D$ annihilate every monomial
with an odd coordinate exponent. Cyclic symmetry and $\phi^2=\phi+1$
give
\[
 \langle x_i^2\rangle_D=\frac13,\qquad
 \langle x_i^4\rangle_D=\frac{1+\phi^4}{3(\phi+2)^2}=\frac15,\qquad
 \langle x_i^2x_j^2\rangle_D=\frac{\phi^2}{3(\phi+2)^2}=\frac1{15}
 \quad(i\ne j).
\]
Together with the constant moment these exhaust the monomials of degree
at most five and equal the spherical moments. Thus $D$ is a $5$-design.
The matrix $T$ is orthogonal of determinant one, so rotation preserves
all design identities and antipodality.

It remains to compare $a,T^ka'$ for $a,a'\in D$ and $k=1,2$.
The rotation angle $\vartheta$ satisfies
$\cos\vartheta=99/101$, $\sin\vartheta=20/101$.
For $a=a'$,
\[
 1-a\cdot T^ka=(1-\cos(k\vartheta))(a_1^2+a_2^2).
\]
The smallest possible $a_1^2+a_2^2$ is $1/(\phi+2)$.
Moreover $1-\cos\vartheta=2/101$ and
$1-\cos(2\vartheta)=800/10201\ge2/101$.
The minimum in this case is exactly $b$.
For distinct vertices of $D$, direct substitution gives
$a\cdot a'\in\{-1,\pm1/\sqrt5\}$, so $|a-a'|>1$.
For $k=1,2$, $\|T^k-I\|\le40/101<2/5$. Hence
\[
 |a-T^ka'|\ge |a-a'|-|(T^k-I)a'|>3/5,
 \qquad 1-a\cdot T^ka'>9/50>b.
\]
This proves the lower bound for every pair. Replacing $a'$ by $-a'$,
which is again in $D$, proves the upper bound $2-b$.
Taking $a=(0,1,\phi)/r$ and $k=1$ attains the lower bound.
\end{proof}

\begin{lemma}[A proper three-colouring]\label{l:colour}
Every $\Z_L\times\Z_N$ with $L,N\ge3$ has a proper three-colouring.
\end{lemma}
\begin{proof}
On the representatives $0,\ldots,m-1$ of $\Z_m$ define
\[
 a_m(i)=\begin{cases}2,&m\text{ odd and }i=m-1,\\
                         i\bmod2,&\text{otherwise}.
          \end{cases}
\]
Adjacent vertices of the cycle have different colours: the usual
alternation works when $m$ is even; when $m$ is odd the last vertex has
colour two and its two neighbours have colours zero and one. Then
\[
 \kappa(i,j)=a_L(i)+a_N(j)\pmod3
\]
is proper on the torus, since an edge changes exactly one summand by a
nonzero residue. This includes $m=3$.
\end{proof}

Assign $D_x=D^{(\kappa(x))}$ and $w_x\equiv1/12$, and set
$\Omega=\prod_xD_x$, $\pi(s)=\prod_x5w_x(s_x)$.
The coherent-state frame map $\iota$ is the same tensor evaluation map
as in the manuscript and is isometric by the tight-frame lemma. Its
ground-state measure is
\begin{equation}\label{e:mucirc}
 \mu^{\af}(s)=\frac{\pi(s)|\iota B(s)|^2}{\|B\|^2}
 =\frac1{Z^{\af}}\prod_x\frac1{12}
       \prod_{\{x,y\}\in E}(1-s_x\cdot s_y).
\end{equation}
Here $|q_{xy}(\zeta(a),\zeta(a'))|^2=(1-a\cdot a')/2$ as before.
Every factor is at least $b$, so $\mu^{\af}$ has full support and
$\iota B$ never vanishes. Bounds in \eqref{e:separation} apply to
individual bond factors; products have the corresponding powers.

\subsection{The frame comparison, including orientation signs}

\begin{theorem}[Replacement for Theorem QD]\label{t:QD}
For every $L,N\ge3$ and every proper colouring used above,
\begin{equation}\label{e:QDall}
 H_{L,N}\big|_{B^\perp}\ \ge\ \frac{g_\star}{2}
                  \lambda(\mu^{\af}) I.
\end{equation}
Thus $\ker H_{L,N}=\C B$ and
$\gap(H_{L,N})\ge(g_\star/2)\lambda(\mu^{\af})$.
More generally the same assertion holds on every nonempty finite simple
$4$-regular graph equipped with positive $4$-design frames whose
adjacent alphabets are disjoint.
\end{theorem}
\begin{proof}
Let $H_x=\sum_{e\ni x}P_e$, $\mathcal S_x=\ker H_x$.
There are four distinct neighbours, so $H_x$ is the same star operator
on five tensor factors at every $x$, irrespective of edges between its
neighbours. Its smallest positive eigenvalue $g_\star$ is positive and
independent of the graph. The operator is nonzero, for instance on the
aligned highest-weight vector. Thus
\[
 H_x\ge g_\star(1-\Pi_{\mathcal S_x}),\qquad \sum_xH_x=2H.
\]
For fixed $x$ put
$\beta_x(t;\sigma)=\prod_{y\sim x}q_{xy}(\zeta(t),\zeta(\sigma_y))$.
Every factor is nonzero. Write
\begin{equation}\label{e:cx}
 q_e=\epsilon_{e,x}q_{xy}\quad(e=\{x,y\}),\qquad
 c_x=\prod_{e\ni x}\epsilon_{e,x}\in\{\pm1\},\qquad
 \rho_x=c_x\prod_{e\not\ni x}q_e\bigl((\zeta(s_z))_{z\in e}\bigr).
\end{equation}
Then $\iota B(s)=\beta_x(s_x;s_{N(x)})\rho_x(s_{\Lambda\setminus x})$.
The local projection $P_x$ in $\ell^2(\Omega,\pi)$ is the fibrewise
projection onto $\C\beta_x$ in $\ell^2(D_x,5w_x)$.
The map
\[
 U:L^2(\mu^{\af})\longrightarrow\ell^2(\Omega,\pi),\qquad
 UF=F\,\iota B/\|B\|
\]
is unitary, because its multiplier is nowhere zero and its squared
modulus gives exactly \eqref{e:mucirc}. Conditional probabilities at $x$
are proportional to $w_x(t)|\beta_x(t;\sigma)|^2$, so the fibrewise
projection formula gives $UE_x=P_xU$.

If $\psi\in\mathcal S_x$, divisibility by the four incident determinants
gives $\psi=(\prod_{e\ni x}q_e)\chi$, where $\chi$ is independent of
$z_x$. Therefore
\[
 \iota\psi(t,s_{\Lambda\setminus x})
   =c_x\beta_x(t;s_{N(x)})\chi((\zeta(s_y))_{y\ne x})\in\Ran P_x.
\]
The isometry of $\iota$ implies
$\|(1-P_x)\iota\varphi\|\le
  \|(1-\Pi_{\mathcal S_x})\varphi\|$.
For $\varphi\perp B$, let $F=U^{-1}\iota\varphi$. Then
$\mu^{\af}(F)=0$ and $\Var_{\mu^{\af}}F=\|\varphi\|^2$.
Consequently
\[
 \lambda(\mu^{\af})\|\varphi\|^2
 \le\sum_x\|(1-P_x)\iota\varphi\|^2
 \le\sum_x\langle\varphi,(1-\Pi_{\mathcal S_x})\varphi\rangle
 \le\frac2{g_\star}\langle\varphi,H\varphi\rangle.
\]
This proves \eqref{e:QDall}. Full support makes the finite single-site
chain irreducible: coordinates can be changed successively with positive
transition probabilities. Its gap is positive, giving uniqueness again.
Every step in this proof applies to the general graph in the statement.
\end{proof}

\subsection{The replacement for the rotated convention}

The torus measure is always \eqref{e:mucirc}; remove the global definition
$n_x=(-1)^{x_1+x_2}s_x$ on the torus. On the plane, however, the sign
$\varepsilon_z=(-1)^{z_1+z_2}$ is well defined. For any proper colouring
$\kappa:\Z^2\to\{0,1,2\}$ define the ferromagnetic specification
\begin{equation}\label{e:muDW}
 \mu^{\ff,\eta}_{\kappa,W}(n_W)
 =\frac1{Z^{\ff,\eta}_{\kappa,W}}\prod_{z\in W}\frac1{12}
   \prod_{\substack{\{z,z'\}\subset W\\z\sim z'}}(1+n_z\cdot n_{z'})
   \prod_{\substack{z\in W,z'\notin W\\z\sim z'}}(1+n_z\cdot\eta_{z'}),
 \quad n_z\in D^{(\kappa(z))}.
\end{equation}
Antipodality ensures that $s_z\mapsto\varepsilon_zs_z$ preserves every
alphabet and its uniform weights. It changes a plane antiferromagnetic
law into \eqref{e:muDW}, with boundary letters
$\eta_{z'}=\varepsilon_{z'}s_{z'}$, and intertwines every single-site
conditional expectation. Hence it preserves the heat-bath gap.
This statement is local to the plane and does not require a periodic sign.

\section{Section 3: the finite-size criterion is unchanged}
\label{s:knabe}

Use manuscript Theorem \texttt{t:K} for the positive measure
$\mu^{\af}$ of \eqref{e:mucirc}. Its hypotheses already allow
site-dependent alphabets, arbitrary positive nearest-neighbour weights,
and all integers $L,N\ge\ell+1$, $\ell\ge2$. In particular, writing
$g(\ell)$ for the minimum of all conditional $\ell\times\ell$ box gaps,
\begin{equation}\label{e:K}
 \lambda(\mu^{\af})\ge\frac{\ell}{\ell-1}
                      \left(g(\ell)-\frac1\ell\right).
\end{equation}
The weighted version and its constants $t(\ell),\kappa_\ell$ also retain
exactly their statements. Neither proof uses global bipartiteness.
The parity of superlevel interval lengths in the autocorrelation lemma
is a property of symmetric intervals on $\Z$, independent of $L,N$.
Retain the ADJ-MINOR corrections: nontrivial site alphabets, and the
separate zero-total-Dirichlet-form case before division in the weighted
Cauchy--Schwarz step. Our alphabets all have twelve letters.

\begin{remark}[Replacement for Remark \texttt{r:Kpins}]\label{r:Kpins}
The minimum in $g(\ell)$ runs only over outside configurations that
occur on the torus. If a period is $\ell+1$, opposite exterior pins
can represent the same torus spin. After lifting and making the local
gauge they can have opposite signs. Lemma~\ref{l:boxcell} realizes
them as admissible, separately specified plane boundary letters.
The uniform bound for all such boundary conditions therefore covers
this case without requiring independent boundary spins on the torus.
\end{remark}

\section{Replacement for Section 4: inhomogeneous cells and box gaps}
\label{s:mixing}

\subsection{Cells and uniform finite energy}

Fix any proper colouring $\kappa:\Z^2\to\{0,1,2\}$. For any
$a\in\Z^2$ use the cells
\[
 C_c=a+2c+\{0,1\}^2,\qquad c\in\Z^2.
\]
The cell alphabet is $S_c=\prod_{z\in C_c}D^{(\kappa(z))}$; it has
$12^4=20736$ elements. There are four internal bonds per cell and eight
bonds crossing its boundary. Formula \eqref{e:muDW} becomes a strictly
positive nearest-neighbour cell specification with site-dependent
single-cell and pair potentials. It has range one in the cell
$\ell^\infty$ metric. No invariance of the colouring is assumed.

Fix an enumeration of the twelve vertices of each frame. This identifies
every $S_c$ with the same abstract alphabet $\{1,\ldots,12\}^{\{0,1\}^2}$,
of size $20736$. All inhomogeneity can therefore be placed in the
potentials, even if a formulation of a finite-state lemma uses a single
alphabet. Equality and Hamming distance compare letters at the same
site, so these identifications preserve them.

Put
\begin{equation}\label{e:energy}
 r_b=\frac{b}{2-b},\qquad
 \beta=\frac1{12}r_b^4,\qquad
 \alpha=\frac1{12^4}r_b^{12},\qquad c_0=\frac{\beta^5}{4}.
\end{equation}
Every physical single-site conditional letter probability is at least
$\beta$; every conditional cell letter probability is at least $\alpha$.
Indeed, for a physical letter at most four factors depend on it, and
each is in $[b,2-b]$. For a cell, at most twelve factors depend on its
configuration and all $12^4$ letters have the same product site weight.
Dividing the least possible weight by the number of letters times the
largest possible weight proves the bounds. The same bounds hold in
every finite-volume conditional law. Conditioning on fewer spins gives
a weighted average of fully specified one-site conditional probabilities,
so a prescribed assignment on $k$ cells has conditional probability at
least $\alpha^k$, by the chain rule.

These finite-energy bounds concern admissible frame-valued boundary
letters. We do not assert them for arbitrary unit-vector boundary pins;
such pins occur only in the separate continuous estimates.

\begin{lemma}[Replacement for Lemma 4.2]\label{l:compare}
On any finite union of these cells, under every admissible boundary
condition, the physical single-site gap is at least $c_0$ times the
cell heat-bath gap, uniformly in $\kappa$ and $a$.
\end{lemma}
\begin{proof}
We use the following elementary bound, valid for different finite
alphabets at different sites. On $k\ge1$ sites, if all one-site conditionals
are at least $\beta$, the single-site gap is at least
$\beta^{k+1}/k$. Here is a direct proof. The chain rule gives every
configuration probability at least $\beta^k$. In the conditional-variance
expansion the coefficient of every directed single-coordinate edge is
at least $\beta^{k+1}/2$. Join any ordered pair of configurations by
changing their coordinates in a fixed order. The path has length at
most $k$, so Cauchy--Schwarz bounds its squared increment by $k$ times
the sum of squared edge increments. In the pair expansion of the
variance, the total probability of pairs using any given directed edge
is at most one. Comparison of these two finite sums gives
$\Var f\le k\beta^{-k-1}\sum_x\mu(\Var_xf)$.

Condition on the exterior of one four-site cell. Its physical gap is
at least $\beta^5/4=c_0$, so
\[
 \mu(\Var_C f)\le c_0^{-1}\sum_{x\in C}\mu(\Var_x f).
\]
Sum over the disjoint cells and apply the cell Poincar\'e inequality.
Every physical site is counted once, proving the comparison without
an additional factor of four.
\end{proof}

\begin{lemma}[Replacement for Lemma 4.1, including equality at the threshold]
\label{l:boxcell}
Let $\ell=2t\ge2$, $L,N\ge\ell+1$, and let $Q$ be any translated
$\ell\times\ell$ box in the torus. For any configuration $\sigma$
outside $Q$, the conditional antiferromagnetic law
$\mu^{\af}(\,\cdot\mid\sigma)$ is carried, by a sitewise bijection,
to a law \eqref{e:muDW} on a plane box made of $t\times t$ cells,
for a proper plane colouring and an admissible cell boundary condition.
The bijection intertwines all single-site conditional expectations.
\end{lemma}
\begin{proof}
Choose an integer lift $Q'\subset\Z^2$ and let
$p:\Z^2\to\Z_L\times\Z_N$ be the quotient map. It is injective on
$Q'$. It identifies exactly the internal nearest-neighbour bonds of
$Q'$ with the torus bonds whose endpoints lie in $Q$. To check the
second assertion, coordinate differences in $Q'$ have absolute value
at most $\ell-1\le L-2,N-2$ as appropriate. A difference congruent to
$\pm1$ modulo a period can thus be only $\pm1$; a difference congruent
to zero is zero. No extra wrap-around internal bond is possible.

The lift $\widetilde\kappa(z)=\kappa(p(z))$ is a proper colouring on
the whole plane. Let $\varepsilon_z=(-1)^{z_1+z_2}$ and set
\begin{equation}\label{e:localgauge}
 n_z=\varepsilon_zs_{p(z)}\quad(z\in Q'),\qquad
 \eta_{z'}=\varepsilon_{z'}\sigma_{p(z')}
     \quad(z'\notin Q',\ z'\sim Q').
\end{equation}
Every exterior site in this formula maps outside $Q$. By antipodality,
all new letters belong to their prescribed frames. On an internal or
boundary plane edge, $\varepsilon_z\varepsilon_{z'}=-1$, so each
antiferromagnetic factor becomes exactly the corresponding positive
factor $1+n_z\cdot n_{z'}$ or $1+n_z\cdot\eta_{z'}$. Site weights are
unchanged. This gives precisely \eqref{e:muDW} after normalization.

If $L=\ell+1$ or $N=\ell+1$, different exterior plane sites may have
the same torus image. They still receive separate, well-defined letters
by \eqref{e:localgauge}. In an odd period those letters can differ by
a sign. A boundary condition on the plane imposes no relation between
distinct exterior sites. Each exterior nearest-neighbour site of a
rectangle touches exactly one interior site; the two torus boundary
bonds therefore remain two separate factors, even if they share an
exterior torus endpoint. This proves the identification in these cases
as well. Extend $\eta$ arbitrarily on the unused exterior sites.

Choose the cell origin to be the lower left corner of $Q'$. Then $Q'$
is exactly a $t\times t$ cell box. Finally, any coordinatewise bijection
of a finite product that transports a probability law transports its
conditional expectation at that coordinate, by the finite conditional
probability formula. This proves the last assertion and preservation
of the heat-bath gap.
\end{proof}

\subsection{The uniform window hypothesis and its stopped coupling}

Here all geometry uses cell coordinates. For a translate
$C_u=a+\{-u,\ldots,u\}^2$, with integer $u\ge1$, and $y\notin C_u$
with $\dist(y,C_u)=1$, write
\begin{align*}
 s&=u^{1/4},\quad k_i=\lfloor is\rfloor,\quad i_1=\lfloor s\rfloor+1,\\
 F_i&=\{c\in C_u:|c-y|_\infty>(i-1)s\},\\
 D_i&=\{c\in C_u:|c-y|_\infty=k_i\},\\
 A_i&=\{c\in D_i:\dist(c,\Z^2\setminus C_u)\le l_0\},
 \qquad B_i=D_i\setminus A_i.
\end{align*}
Let $I_0\subset\{2,\ldots,i_1\}$ be the indices for which $D_i$ is
at distance less than $s/3$ from a corner of $C_u$.

\begin{definition}[Replacement for Definition \texttt{d:W}]\label{d:W}
A family of specifications has the uniform window estimate if there
are an integer $l_0\ge1$, $p_0\in(0,1]$ and an integer $M_\star\ge1$
such that for every specification in the family, every translated
$C_u$ with integer $u\ge M_\star$, every $y$ as above, every integer
$2\le i\le s$ outside $I_0$, and every pair of admissible boundary
conditions on $F_i$,
\begin{equation}\label{e:W}
 \TV(\widehat\mu_{F_i,B_i}^{\eta},
       \widehat\mu_{F_i,B_i}^{\xi})\le1-p_0.
\end{equation}
\end{definition}

The manuscript shell geometry (MOSPROOF Lemma 4.2) gives the following
whenever $u\ge M_{\rm g}(l_0)$, where
\[
 M_{\rm g}(l_0)=\max\{16,6(l_0+1)\}^4:
\]
\begin{equation}\label{e:shellfacts}
 |I_0|\le1,\qquad |A_i|\le8l_0\ (i\notin I_0),\qquad D_i\subset F_i.
\end{equation}
Also, the neighbours of $F_i$ in $C_u\setminus F_i$ lie on $D_{i-1}$,
and $y$ does not neighbour $F_i$ for $i\ge2$. These are geometric
statements, valid in every translate and independent of the alphabets
or the potential. Reflections used in their proof transform coordinates
only; they require no reflection invariance of a specification.

\begin{lemma}[Stopped surgery, with no stationarity hypothesis]
\label{l:stopped}
Suppose a positive, range-one finite-state specification has one-site
conditional probabilities at least $\alpha$ and satisfies
Definition~\ref{d:W} at every translate. Put
$\delta=p_0\alpha^{16l_0}$. If
$u\ge\max\{M_\star,M_{\rm g}(l_0)\}$ and two boundary conditions of
$C_u$ agree off one exterior site $y$, then for every integer
$4\le h\le\lfloor u^{1/4}\rfloor$ their full-volume laws have a
coupling such that
\begin{equation}\label{e:stopped}
 \Pr(n_{Q_h}\ne n'_{Q_h})\le(1-\delta)^{h-3},\qquad
 Q_h=\{c\in C_u:|c-y|_\infty>hu^{1/4}\}.
\end{equation}
The assertion allows site-dependent finite alphabets and potentials.
\end{lemma}
\begin{proof}
If $\dist(y,C_u)>1$, the two laws coincide by the range-one Markov
property. Otherwise begin with any coupling, for example the product
coupling, and resample $F_i$ successively for $i=2,\ldots,h+1$,
skipping indices in $I_0$. Conditional resampling leaves each marginal
unchanged: on each exterior fibre its new law is precisely its old
conditional law. This is a finite sum identity, independent of any
translation invariance.

If the input pair agrees on $D_{i-1}$, the two conditional laws of
$F_i$ agree, by \eqref{e:shellfacts} and equality of the fixed exterior
boundary off $y$. Use the same sample in both components. Otherwise
maximally couple the two $B_i$ marginals, complete each sample by its
conditional law, and independently resample $A_i$ in both components
given their respective complements. For $i\le h$, \eqref{e:W} gives
agreement on $B_i$ with probability at least $p_0$. For any fixed
assignment on $A_i$, both independent resamplings choose it with
probability at least $\alpha^{2|A_i|}$. Thus agreement on $D_i$ has
conditional probability at least $\delta$ at each such step. Each
operation preserves both required marginals. At the final step $h+1$
we need no new success bound if $h+1>s$.

Let $J$ be the indices $i\in\{2,\ldots,h\}$ with both $i$ and $i+1$
outside $I_0$. There are at least $h-3$ of them. Success at any such
$i$ forces identical resampling on $F_{i+1}$ at the next step.
Agreement there persists at later steps: they change only smaller
$F_j$, and their separators $D_{j-1}$ are contained in $F_{i+1}$.
For $j\ge i+2$ the latter inclusion follows from
$\lfloor(j-1)s\rfloor>is$, since $s\ge16$. It gives final agreement
on $Q_h=F_{h+1}$. By conditioning successively on the history at the
trial indices, the probability that all trials fail is at most
$(1-\delta)^{|J|}\le(1-\delta)^{h-3}$. Independence between trials is
not asserted. All couplings exist on finite sets. No operation compares
letters at different sites.
\end{proof}

Taking $h=\lfloor s\rfloor$ also proves the original conclusion of
Proposition \texttt{p:surgery}, with exponent $s-4$, for every translate
and every specification in the family. The terminal step is a
propagation step; its success probability is not inferred from a window
estimate at an index greater than $s$.

\subsection{A box-gap theorem that consumes the uniform window estimate}

\begin{theorem}[Replacement for Corollary \texttt{c:gapcell}]
\label{c:gapcell}
Let a family of positive range-one finite-state specifications have a
common conditional lower bound $\alpha\in(0,1]$ and the uniform window
estimate $(l_0,p_0,M_\star)$. Neither a common spatial pattern of
alphabets nor translation invariance is required. Set
\[
 d=\min\{1/2,p_0\alpha^{16l_0}\}.
\]
Choose an integer $n$ such that
\begin{equation}\label{e:nchoice}
 n\ge44,\quad n^8\ge M_\star,\quad
 n^2\ge\max\{50000,6(l_0+1)\},\quad
 n^{12}d^{20}\ge50000\,2^{20}\,20!.
\end{equation}
Put
\begin{equation}\label{e:cellconstants}
 M=n^8,\quad L_{\rm cell}=4M+1,\quad v=L_{\rm cell}^2,\quad
 m_0=\frac{\alpha^{v+1}}{1000v}>0.
\end{equation}
Every translated rectangle with both cell sides at least $L_{\rm cell}$
has single-cell heat-bath gap at least $m_0$, uniformly over the family
and over all admissible boundary conditions.
\end{theorem}
\begin{proof}
The choices in \eqref{e:nchoice} exist by taking $n$ sufficiently large.
For each $u\in[M,2M]$ apply Lemma~\ref{l:stopped} with $h=n$ to every
translated square $C_u$. Its hypotheses hold because
$u^{1/4}\ge n^2\ge\max\{16,6(l_0+1)\}$ and $u\ge M_\star$.
The exceptional near region $C_u\setminus Q_n$ has at most $25n^6$
cells, since $nu^{1/4}\le2n^3$. The whole square has at most $25n^{16}$
cells. In the actual coupling of that lemma, the expected Hamming
distance is consequently at most
\begin{equation}\label{e:Kbound}
 K_n=25n^6+25n^{16}(1-d)^{n-3}\le M/1000.
\end{equation}
For the last inequality, the finite binomial expansion gives
\[
 (1-d)^{n-3}\le\frac{2^{20}20!}{d^{20}n^{20}}\quad(n\ge44):
 \quad (1-d)^{-(n-3)}\ge\binom{n-3}{20}d^{20}
                         \ge\frac{(n/2)^{20}}{20!}d^{20}.
\]
The first term in \eqref{e:Kbound} is at most $n^8/2000$ by
$n^2\ge50000$, and the second is at most $n^8/2000$ by the final
inequality of \eqref{e:nchoice}. This supplies the square Hamming
coupling criterion at every translate, directly, without transporting
it by an invariance assumption.

Here is the finite-square gap step, included to specify precisely the
part of MOSPROOF being used. In a rectangle
$W=\{1,\ldots,N_1\}\times\{1,\ldots,N_2\}$ with both sides at least
$L=4M+1$, use labels
\[
 \mathcal A=\{M,\ldots,2M\}\times\{1-L,\ldots,N_1\}
                        \times\{1-L,\ldots,N_2\}.
\]
For a label $(u,a_1,a_2)$ take a square $B_a$ of side $l=2u+1$ with
lower left coordinates
$b_j=\min\{\max\{a_j,1\},N_j-l+1\}$. Retain multiplicities.
The purely geometric counting lemma MOSPROOF Lemma 6.1 gives, for
each $x\in W$,
\begin{equation}\label{e:counts}
 4M^3\le T_x:=\#\{a:x\in B_a\}\le1000M^3,\qquad
 S_x:=\#\{a:x\notin B_a,\dist(x,B_a)=1\}\le1000M^2.
\end{equation}
These counts do not refer to the Gibbs specification. In particular
they do not use translation invariance. One way to see their uniformity
at the boundary is that clamping preserves containment for the $l$
nominal starts in $[x_j-l+1,x_j]$. This gives the lower bound.
An interval of starts containing $x_j$ has at most $3L$ preimages,
giving $T_x\le(M+1)(3L)^2$. For $S_x$, overcount by a directed side:
the start in one coordinate must be $x_j+1$ or $x_j-l$. A strictly
interior such start has one preimage, and the other coordinate has at
most $4L$ preimages, giving $16(M+1)L$. Endpoint starts can occur for
at most one $u$ on each of the four directed sides, have at most $2L$
preimages in the constrained coordinate, and contribute at most
$32L^2$. With $L\le5M$ and $M+1\le2M$ these are exactly
\eqref{e:counts}.

Let $A=|\mathcal A|$ and $P=A^{-1}\sum_a E_{B_a}$. Couple its output
laws for two configurations differing just at $x$. If $x\in B_a$,
use the same block sample, erasing the disagreement. If
$\dist(x,B_a)>1$, use identical samples and retain just that
disagreement. In the other $S_x$ cases, \eqref{e:Kbound} bounds the
expected number of new disagreements in the block by $K_n$. Thus
\[
 \E d_H\le1-\frac{T_x-K_nS_x}{A}\le1-\frac{M^3}{A}.
\]
The Hamming Lipschitz seminorm consequently contracts by
$1-M^3/A$ under $P$. Each $E_{B_a}$ is an orthogonal projection in
$L^2(\mu_W)$, so $P$ is self-adjoint. Apply the Lipschitz bound to
each nonconstant real eigenfunction; its seminorm is positive on the
connected product space. Every such eigenvalue of $P$ has absolute
value at most $1-M^3/A$. Hence the unnormalized block generator
$\sum_a(I-E_{B_a})=A(I-P)$ has gap at least $M^3$.

Each conditional block law has at most $v$ cells and its one-cell
conditionals remain at least $\alpha$. The elementary small-volume
bound proved in Lemma~\ref{l:compare}, with $k\le v$ and $\beta$
replaced by $\alpha$, gives its single-cell gap at least
$g=\alpha^{v+1}/v$. Integrating its Poincar\'e inequality over each
exterior fibre and summing over blocks gives
\[
 M^3\Var_\mu f
 \le\sum_a\mu(\Var_{B_a}f)
 \le g^{-1}\sum_x T_x\mu(\Var_xf)
 \le\frac{1000M^3}{g}\sum_x\mu(\Var_xf).
\]
Thus the single-cell gap is at least $g/1000=m_0$.
This is MOSPROOF Theorem 6.3 with all its hypotheses verified. The
preceding stopped argument is its window-to-coupling step with the
translation-invariance sentence replaced by the all-translates
quantifier. All sums, fibres, couplings and eigenvectors in the proof
are finite. Bounds for complex functions follow by adding the real
and imaginary Poincar\'e inequalities.
\end{proof}

Theorems \texttt{t:MOS21}, \texttt{t:MO41}, Remark \texttt{r:finiteS},
and the compatible-set scales are removed from the proof chain.
Their original translation-invariant formulations are not applied to
the inhomogeneous specification. The new conclusion is a gap on all
sufficiently large rectangles; no claim on arbitrary compatible unions
is needed here.

\subsection{Replacement proof of the main theorem from the window estimate}

\begin{theorem}[Replacement for Theorem \texttt{t:mainW}]
\label{t:mainW}
Assume the uniform window estimate for the ferromagnetic cell models
of every proper plane three-colouring, with the same parameters.
Let $m_0,L_{\rm cell}$ be given by \eqref{e:cellconstants} using
\eqref{e:energy}. Choose an integer
\begin{equation}\label{e:ellchoice}
 t\ge\max\{L_{\rm cell},\lceil1/(c_0m_0)\rceil,1\},\qquad
 \ell=2t,\qquad L_0=\ell+1,\qquad c=g_\star c_0m_0/4.
\end{equation}
For every $L,N\ge L_0$, $\ker H_{L,N}=\C B$ and
$\gap(H_{L,N})\ge c$. The same constants work for every proper
three-colouring of every such torus used to construct its frame measure.
\end{theorem}
\begin{proof}
For every placement of an $\ell$-box and every exterior configuration,
Lemma~\ref{l:boxcell} identifies its conditional law, up to the local
gauge, with a $t\times t$ cell law of a proper plane colouring.
Theorem~\ref{c:gapcell} gives its cell gap at least $m_0$, and
Lemma~\ref{l:compare} gives its physical gap at least $c_0m_0$.
Consequently $g(\ell)\ge c_0m_0$ and $1/\ell\le c_0m_0/2$.
The finite-size criterion \eqref{e:K} now gives
\begin{equation}\label{e:classicalgap}
 \lambda(\mu^{\af})\ge\frac{\ell}{\ell-1}
                 \left(c_0m_0-\frac1\ell\right)\ge c_0m_0/2.
\end{equation}
Theorem~\ref{t:QD} proves the stated quantum gap, with the constant in
\eqref{e:ellchoice}. Uniqueness also follows independently from
Proposition~\ref{p:graphkernel}. Every constant was chosen before the
periods, the torus colouring, the box placement and its boundary data.
\end{proof}

\section{Section 5: uniformity of the cut and window estimates}
\label{s:cut}

\subsection{Quadrature for different frames}

\begin{lemma}[Replacement for the exactness lemma and its following paragraph]
\label{l:exact}
Assign any positive weighted $4$-design to each site of a finite set.
A product of factors $1\pm n_x\cdot n_y$ and $1\pm n_x\cdot v$,
with fixed unit pins $v$, has the same frame sum as its spherical
integral whenever each variable occurs in at most four factors.
The assertion permits a different frame and a different sign choice
at every site or factor.
\end{lemma}
\begin{proof}
The integrand is a polynomial of separate degree at most four.
Replace each frame sum by the spherical integral at that site in a
fixed order. The design identity applies with all other variables
held fixed. Summation or integration in other variables cannot raise
the degree of the variable being treated. Finite Fubini, or ordinary
Fubini for a bounded polynomial on a compact product, completes the
proof. The application here has distinct endpoints on every bond;
if a self-factor were allowed in the general wording, it is first
replaced by its constant value on unit spins.
\end{proof}

In particular every partition function in the cut construction for
\eqref{e:muDW} equals its continuous counterpart, independently of
the proper colouring. The pins can be arbitrary unit vectors in
this equality. This does not assert equality of window measures for
different frames: the sample spaces and the window's own weight
can differ. Only the summed-out partition functions are transferred.

\subsection{The cut ratio with arbitrary site frames}

We retain the continuous ferromagnetic model, its original legs and
removed-site convention, and Hypothesis (P) of the manuscript. Explicitly,
for every finite $G\subset\Z^2$, every unit-vector pin assignment $\tau$
on its original legs, every $Y\subset G\setminus S(G)$, where
$S(G)=\{x\in G:d(x,\Z^2\setminus G)\le2\}$, and every pair
$u,v\in G\setminus Y$ at distance two with a common neighbour in $Y$,
\begin{equation}\label{e:P}
 \left|\E^\tau_{G\setminus Y}(n_u\cdot n_v)
             -\E^0_{G\setminus Y}(n_u\cdot n_v)\right|
 \le A\theta^{\min\{d(u,S(G)),d(v,S(G))\}},\qquad
 A\ge2,\quad0<\theta<1.
\end{equation}
All expectations in this hypothesis are continuous; there are no
torus identifications.

The cut-ratio lemma of the manuscript holds for arbitrary site frames
in \eqref{e:muDW}. Here is the full verification of the additional
uniformity being used. Let $F$ be a finite physical domain and
$B\subset F$ a window not adjacent to its exterior. Put $G=F\setminus B$.
For a window letter assignment $s$, define
\[
 W_B(s)=\prod_{x\in B}w_x(s_x)
                \prod_{\substack{\{x,y\}\subset B\\x\sim y}}(1+s_x\cdot s_y).
\]
With pins $\eta s$ on the original legs of $G$, quadrature gives
\begin{equation}\label{e:windowlaw}
 \mu_F^\eta(n_B=s)=\frac{W_B(s)Z_G^{\eta s}}{Z_F^\eta}.
\end{equation}
Let a cut $X\subset G$ separate $G\setminus X=G_1\sqcup G_2$, with
no bonds between the pieces, no exterior pins on $G_1$ and no window
pins on $G_2$. Write $\Psi_D=Z_D^\tau/Z_D^0$ and
\[
 e(\eta,s)=\log\Psi_{G\setminus X}(\eta s)-\log\Psi_G(\eta s).
\]
Factorization gives
$Z_{G\setminus X}^{\eta s}=Z_{G_1}^{s}Z_{G_2}^{\eta}$ and
$Z_{G\setminus X}^{0}=Z_{G_1}^{0}Z_{G_2}^{0}$. Therefore
\begin{equation}\label{e:nu}
 \mu_F^\eta(n_B=s)=c(\eta)\nu(s)e^{-e(\eta,s)},\qquad
 \nu(s)\ \text{proportional to}\ W_B(s)Z_{G_1}^{s}.
\end{equation}
The law $\nu$ depends on the chosen frames and the window, but not
on $\eta$.

All continuous normalizers here are positive, also for arbitrary
pins: choose a unit vector not antipodal to any of the finitely
many pins and put every free spin in a sufficiently small neighbourhood
of it. All factors are strictly positive on this open product set.
The empty-domain partition function is one. For the actual frame
model, $W_B(s)>0$ by \eqref{e:separation}, so $\nu(s)>0$.
If $|e(\eta,s)|\le\epsilon$ uniformly, summing \eqref{e:nu} shows
$e^{-\epsilon}\le c(\eta)\le e^\epsilon$. Thus
\begin{equation}\label{e:cutratio}
 e^{-2\epsilon}\le\frac{\mu_F^\eta(n_B=s)}{\nu(s)}
                   \le e^{2\epsilon}.
\end{equation}
For two boundary conditions the density ratio is bounded in both
directions by $\rho=e^{4\epsilon}$. If $A$ is the set where the first
law exceeds the second and $x$ is the second law's mass of $A$,
their TV distance is at most
$\min\{(\rho-1)x,(\rho-1)(1-x)/\rho\}\le(\rho-1)/(\rho+1)$.
It is therefore at most $\tanh(2\epsilon)$. This proof involves
the actual window weights and does not require the window itself
to have a frame-independent law.

\begin{theorem}[Replacement for Theorem \texttt{t:cutMOS}]
\label{t:cutMOS}
Assume \eqref{e:P}. The cell models \eqref{e:muDW} for every proper
plane three-colouring and every cell origin have the uniform window
estimate of Definition~\ref{d:W}, with $l_0=20$ and constants
$p_0>0,M_\star$ depending only on $A,\theta$. In particular these
constants are independent of colouring, alphabet cardinality,
translation and torus periods.
\end{theorem}
\begin{proof}
Use exactly the cell sets $U,F_i,B_i$ and physical cut
$X=\partial_\infty U'$ of manuscript Section 5.4. Its geometric
lemmas apply after any translation and any choice of cell origin;
their proofs use only the square grid, not the frames. With the
ADJ-MINOR corrections, $X$ is a nonempty connected finite set,
$b_G(x)=d(x,\Z^2\setminus G)\ge3$ on $X$, and, for $s=u^{1/4}$,
\begin{equation}\label{e:cutsum}
 \sum_{x\in X}\theta^{b_G(x)}
 \le192\theta^5\left(6+\frac1{1-\theta^2}\right)
        +4(2u+1)^2\theta^{s/6-7}.
\end{equation}
The connectedness is used to order the cut sites so that every site
after the first has an earlier neighbour, by a spanning tree. The
nonemptiness follows because $U'$ is finite and nonempty and therefore
has a nonempty exterior vertex boundary.

The unchanged removal lemma \texttt{l:removal}, applied to these sites
and the continuous pins $\eta s$, gives
\begin{equation}\label{e:epsilon}
 |e(\eta,s)|\le0.0707+2.73A+
   1.22A\theta^{-3}\left[
     192\theta^5\left(6+\frac1{1-\theta^2}\right)
           +4(2u+1)^2\theta^{u^{1/4}/6-7}\right].
\end{equation}
Here the decimal constants denote the exact rational upper bounds
in that lemma. Its additive $0.0707$ controls the four-spin term
at the first removal; it is not omitted. Every remaining term is
controlled by \eqref{e:P} in a continuous domain with its original
legs. The lemma therefore applies to every frame-valued pin
assignment encountered here, whatever its site type.

Choose a single integer threshold $M_\star$ so that all the unchanged
geometric requirements hold and the last term inside the brackets
of \eqref{e:epsilon} is at most one for every $u\ge M_\star$.
Such a threshold exists because a fixed polynomial times
$\theta^{u^{1/4}/6}$ tends to zero. Set
\begin{equation}\label{e:windowconstants}
 \epsilon_{\max}=0.0707+2.73A+
  1.22A\theta^{-3}\left[
       192\theta^5\left(6+\frac1{1-\theta^2}\right)+1\right],\qquad
 p_0=\frac2{\exp(4\epsilon_{\max})+1}>0.
\end{equation}
The cut-ratio proof \eqref{e:windowlaw}--\eqref{e:cutratio} applies
with the actual frames at each site, by Lemma~\ref{l:exact}.
It gives \eqref{e:W} with $1-p_0=\tanh(2\epsilon_{\max})$.
All bounds used in this proof precede the choice of colouring or
translate. The physical window configurations and their corresponding
cell letters are in bijection, which preserves TV. This proves the
full all-translates quantifier required in Definition~\ref{d:W}.
\end{proof}

\section{Sections 6--8 and appendices: retained inputs and completion}
\label{s:inputs}

The walk expansion, pair recursion and computer-assisted inputs remain
statements about the continuous ferromagnetic model on finite subsets
of $\Z^2$, with arbitrary unit-vector pins. Keep their adjudicated
versions, with the same constants and local certificate family.
In particular no walk expansion is performed directly on an odd torus.
Here are the precise interfaces needed for the present proof.

In Section 6 use the ADJ-R6 definition of \texttt{d:pattern}: full
kept sets at $\mathrm U$ sites, actual two-bond kept sets at
$\mathrm D$ sites, and the stated special rule at tips. Use the
ADJ-R4 auxiliary networks with positive scalar trace masses; terminal
protected-boundary sites retain their original pin factors. A general
initial two-bond $\mathrm D$ departure has ratio at most $3/4$,
whereas later $\mathrm D$ departures have ratio $3/5$. The exposure
lemma is used only with at most one present initial $\mathrm D$
bond. The common one-point majorant bounds the sum of maxima of the
pinned and free ratios. Impose ADJ-R5's nonreversing initial arrival
bond convention in (H2); the corresponding missing bond is available
at every start used in the recursion.

In Section 7 retain the ADJ-R4 closure lemma, which distinguishes the
live endpoint, the terminal boundary defect and the mass-one underlying
consistent domains. Normalize an initial $\mathrm U$ site's bond to
the removed centre before tracing. These corrections prove the same
recursion, with the same two pair classes and kernel. Its suprema
already range over all finite plane domains and all unit pins, so
they require no change when finite alphabets vary between sites.

In Section 8 and Appendix C retain ADJ-R5's corrected memory-walk
range, coefficient assembly, and nonnegative sub-invariant vector
with its checked zero support. The quotient with periods $(28,28)$
and $(28,-28)$ is a device for bounding trails on the plane, not the
physical torus. A trail between the two distance-two endpoints has
even length because the plane is bipartite; this remains true and
does not restrict $L$ or $N$. The local integer certificates and
closed-form inequalities of Appendix D depend only on arbitrary
unit vectors, not on frames. The total-spin weight parity in that
appendix is representation-theoretic parity, unrelated to torus
periods. Retain ADJ-R6's corrections to certificate provenance and
the qualification of selected numerical minima.

The mod-two even-edge-set argument in Appendix B is likewise entirely
on the plane and is unchanged. The Section 5 geometry uses its
adjudicated nonempty-set and coincident-side details. None of these
inputs asserts translation invariance of a finite frame model.

\begin{proof}[Proof of Theorem~\ref{t:main}]
The corrected computer-assisted Theorem \texttt{t:inputs} and corrected
recursion theorem \texttt{t:rec} give \eqref{e:P}, with some fixed
$A\ge2$ and $0<\theta<1$. Theorem~\ref{t:cutMOS} gives the uniform
window estimate at every translate for every proper plane colouring.
Theorem~\ref{c:gapcell} gives the common constants
$L_{\rm cell},m_0>0$. Lemma~\ref{l:colour} supplies a torus colouring
for every $L,N\ge3$, and Theorem~\ref{t:mainW} proves the uniform
gap with the constants \eqref{e:ellchoice}. Finally,
Proposition~\ref{p:graphkernel} proves uniqueness for every $L,N\ge3$,
without any estimate from Sections 3--8. These dependencies are not
circular: the analytic proof of (P) uses continuous moment identities
and the elementary three-factor bound, not the frame construction
or any finite-volume heat-bath gap.
\end{proof}

\appendix
\section{Additional consequences and small sizes}
\label{s:consequences}

\begin{corollary}[Rectangles, arbitrary colour patterns and classical gaps]
\label{c:extras}
The same constants yield all of the following.
\begin{enumerate}
\item Every proper plane three-colouring, every cell origin, every
rectangle whose physical side lengths are $2t_1,2t_2$, with
$t_1,t_2\ge L_{\rm cell}$, and every admissible boundary condition
have physical single-site ferromagnetic heat-bath gap at least
$c_0m_0$. The same holds in the antiferromagnetic convention after
the local plane gauge.
\item On every torus with $L,N\ge\ell+1$, the antiferromagnetic
frame heat bath of every proper three-colouring has gap at least
$c_0m_0/2$.
\item On an open $r\times s$ square-lattice rectangle, $r,s\ge2$,
the quantum ground-space dimension is
$3^4\,2^{2r+2s-8}$.
\end{enumerate}
\end{corollary}
\begin{proof}
For (1), choose the cell origin at the lower left corner, apply
Theorem~\ref{c:gapcell} to the cell rectangle and then
Lemma~\ref{l:compare}; the local gauge preserves all single-site
operators. Assertion (2) is \eqref{e:classicalgap}, whose proof was
uniform over all proper colourings. For (3), the four corners have
degree two, the other $2r+2s-8$ boundary sites degree three, and
all remaining sites degree four. Apply \eqref{e:graphkernel}.
\end{proof}

There are no exceptions among $L,N\ge3$, including $(3,3)$ and
arbitrarily long tori of width three, for the uniqueness assertion.
No uniform gap for the infinitely many thin tori with
$3\le\min(L,N)<L_0$ is asserted by this argument.

Periods below three are outside the manuscript's model convention:
one must choose whether to suppress loops and repeated bonds or
retain a multigraph. To make the degeneracy issue precise, if one
extends by the simple-graph convention (distinct neighbours only,
loops suppressed), then for $L=2,N\ge3$ every degree is three and
$\dim\ker H=2^{2N}$; for $L=N=2$ it is $3^4$.
For $L=1,N\ge3$ it is $3^N$, for $(L,N)=(1,2)$ or $(2,1)$ it is
$4^2$, and for $(1,1)$ it is $5$. The corresponding swapped cases
are identical. All follow from Proposition~\ref{p:graphkernel}.
In these cases $B_G$ alone has insufficient degree at some sites
to be a vector of the spin-2 tensor product; the residual factors
in \eqref{e:graphkernel} are essential. This is why they are not
additional sizes of the theorem as originally defined.

\section{Replacement constants and source convention}

For the extension, delete the icosahedron--dodecahedron
$\gamma_\pm$ from the chain of required constants and use
\eqref{e:delta}, \eqref{e:energy}, \eqref{e:windowconstants},
\eqref{e:nchoice}, \eqref{e:cellconstants} and \eqref{e:ellchoice},
in that order after $A,\theta$ are available. The auxiliary physical
box side $\ell$ is even because we choose to block it into cells;
this places no parity restriction on the torus. The final torus
threshold can be $\ell+1$, including the boundary-identification
case formerly excluded by parity.

The unchanged constants are $g_\star$, the two Knabe prefactors,
all removal and response coefficients, $c_{\mathrm s},c_{\mathrm t}$,
$C_V,q$, the kernel and its bound on $\rho(M_1)$, and the continuous
$A,\theta$. The new $c,L_0$ need not equal the previous unspecified
values. The gap bound is qualitative; the explicit finite-energy
construction here gives very conservative constants.

This supplement uses these supplied local documents, with their
source labels as stable locators: manuscript Sections 1--8 and
Appendices A--D; MOSPROOF \texttt{proof.tex}, especially Lemmas 3.1,
4.2, 4.3, 6.1 and Theorem 6.3; and the four adjudication reports
ADJ-R4, ADJ-R5, ADJ-R6 and ADJ-MINOR. No MO94/MOS94 mixing theorem
is invoked as an additional black box. The exact construction in
Lemma~\ref{l:threeframes} is proved above without a design citation
or a numerical design test. The independent one-process script
\texttt{check\_exact.py} verifies all its degree-five moments and
cross-frame bounds in $\mathbb Q(\sqrt5)$ and tests finite torus
colourings and boundary-lift cases; these finite tests supplement,
and do not replace, the universal proofs.

\end{document}
